\chapter{Modul Pemeliharaan \& Arsip (Maintenance \& Archive)}

Lompatan tahun ajaran selalu membawa masalah klasik: penumpukan sisa rekam jejak digital yang memberatkan sistem. EduConnect mengatasi ini melalui Modul Pemeliharaan dan Arsip yang didesain untuk mereset ruang kerja secara aman, membersihkan *cache* aplikasi, serta memproteksi data lawas ke dalam balutan arsip terkompresi.

\section{Konsep Dasar \& Matriks Peran}
\begin{itemize}
    \item \textbf{Super Admin}: Satusatunya peran berdaulat yang menduduki puncak takhta wewenang (*god-mode*). Hanya Super Admin yang berhak memencet tombol eksekusi \textit{Archive Data}. Guru, bahkan level Admin sekalipun tidak diperkenankan bermain di area berbahaya ini guna menghambat insiden salah ketik (\textit{fat-finger}) yang menghapus data secara permanen.
\end{itemize}

\section{Skenario Dunia Nyata \& Alur Kerja}
Penghujung bulan Juni 2026, tahun ajaran usai. Siswa kelas 12 telah diluluskan secara seremonial. Fase 1: Super Admin IT sekolah mengakses panel rahasia "Arsip Kelas". Fase 2: Beliau memilih opsi target "Kelas 12 MIPA 1", menandai kotak persetujuan penghapusan memori persisten, dan mengeksekusi "Arsipkan \& Hapus". Fase 3: Di balik layar (secara asinkron), aplikasi mengekstrak riwayat 12 MIPA 1, menyatukannya menjadi \texttt{12-mipa-1\_2026.zip}, memasukkannya ke lumbung \textit{Storage File System}, lalu menebas seluruh data absen dan relasi siswa bersangkutan di *database* tabel utama sehingga kapasitas sistem menjadi lega kembali untuk menyambut Tahun Ajaran 2027.

\begin{figure}[H]
    \centering
    \includegraphics[width=0.9\textwidth]{placeholder_archive_danger_zone.png}
    \caption{Zona Bahaya (Danger Zone) pada halaman Arsip Kelas yang mewajibkan konfirmasi ganda (Double-opt-in).}
    \label{figure:7.1}
\end{figure}

\section{Panduan Penggunaan (Step-by-Step)}

\subsection{Melaksanakan Prosedur Arsip Kelas (Archive Protocol)}
\begin{enumerate}
    \item Pastikan rapor semester/tahunan telah selesai dibagikan, karena prosedur ini akan \textbf{MENGHAPUS DATA UTAMA}.
    \item Pada panel menu, akses wilayah \textbf{Maintenance} lalu klik opsi sekunder \textbf{Arsip Kelas}.
    \item Tampilan antarmuka ini dipisahkan menjadi "Panel Aksi Baru" dan "Tabel Arsip Terdahulu". Berfokuslah pada Panel Aksi (Arsipkan Data Kelas).
    \item Pilih kelas sasaran pada peranti *dropdown* \textbf{Kelas}.
    \item Centang kotak berbunyi \textbf{Saya yakin ingin mengarsipkan dan MENGHAPUS data siswa dari kelas ini beserta seluruh absensinya}. Tanpa centangan ini (\textit{checkbox}), prosedur terkunci demi pengamanan lapis pertama.
    \item Klik tombol warna merah pekat \textbf{Arsipkan \& Hapus Data Kelas}.
    \item Akan timbul peringatan (\textit{alert}) dari *browser*. Konfirmasi kembali bahwa keputusan Anda telah valid (\textit{OK}).
    \item Data siswa dan rekam absensi pada kelas terilih seketika lenyap dari tabel utama \textit{database}.
\end{enumerate}

\subsection{Mengambil Kembali Arsip yang Terkompresi (Download Artifact)}
\begin{enumerate}
    \item Pada halaman yang sama (\textbf{Arsip Kelas}), alihkan pandangan ke bagian bawah halaman pada modul tabel daftar arsip.
    \item Telusuri baris yang memuat catatan nama arsip dengan keterangan kelas yang diretas sebelumnya.
    \item Klik tautan warna biru bertuliskan \textbf{Download Excel} jika Anda butuh \textit{spreadsheet} kilat, atau tautan \textbf{Download ZIP} jika Anda butuh dokumen mentah seluruh data pendamping (termasuk foto, surat sakit/izin lawas, dsb).
\end{enumerate}

\section{Penanganan Masalah (Edge Cases)}
\begin{itemize}
    \item \textbf{Celah Path Traversal}: Beberapa aplikasi abal-abal rentan terhadap penyisipan serangan *Path Traversal* ketika *mengunduh file* arsip (meretas \textit{path} untuk mengunduh `.env` atau *root system*). EduConnect mengunci fungsi `ArchiveController` agar hanya \textit{file} dengan direktori spesifik `/archives/` dan ber-ID valid yang lolos unduh.
    \item \textbf{Timeout Eksekusi ZIP}: Apabila kelas tersebut memilki puluhan ribu rekam absensi beserta foto lampiran izin yang memakan ukuran \textit{Gigabyte}, *server PHP* mungkin akan pingsan (*Time-out limits* 30/60 detik) sebelum berhasil men-ZIP. Disarankan mengeksekusi ini dari sisi \textit{Command Line} (*Artisan Console*) jika ruang tampung datanya tergolong level raksasa.
\end{itemize}
